\subsection{Polar path tetrahedron theorem}
\label{sec:bt1258-polar-path-tetrahedron}

BT1252--BT1255 identify the geometric object underlying the BT1233 diameter-14 word metric.  Given a four-transvection set, form the tetrahedron on its four projective generators.  Mark an edge polar when the two generators have zero symplectic pairing, and nonpolar otherwise.

For the BT1228 representative, the polar edges form a path \(P_4\), and the nonpolar edges form the complementary path \(P_4\).  Thus
\[
K_4=P_4\sqcup P_4.
\]
We call such a configuration a \emph{polar path tetrahedron}.

The full-order regimes have the following edge-graph types:
\[
\begin{array}{c|c|c}
\text{regime} & \text{polar graph} & \text{nonpolar graph}\\
\hline
\operatorname{diam}=10_A & K_2+2I & K_4-e\\
\operatorname{diam}=10_B & 2K_2 & C_4\\
\operatorname{diam}=10_C & 4I & K_4\\
\operatorname{diam}=12 & P_3+I & \mathrm{paw}\\
\operatorname{diam}=14 & P_4 & P_4
\end{array}
\]

Since the diameter-14 full-order regime is a single \(Sp(4,3)\)-orbit, the \(P_4/P_4\) representative globalizes: every diameter-14 full-order four-transvection set is a polar path tetrahedron up to the \(Sp(4,3)\) action.

This explains the balanced closure law
\[
9^3\,24^3
\]
on pairs and
\[
72^2\,648^2
\]
on triples.  The BT1233 fingerprint is therefore the word metric of the unique self-complementary polar/nonpolar edge split.
